\documentclass[12pt,a4paper]{report}

% ---------- Packages ----------
\usepackage[a4paper,left=1.3in,right=0.9in,top=0.85in,bottom=0.85in]{geometry}
\usepackage[T1]{fontenc}
\usepackage{mathptmx}
\usepackage{graphicx}
\usepackage{array}
\usepackage{caption}
\usepackage{float}
\usepackage{xcolor}
\usepackage{titlesec}
\usepackage{fancyhdr}
\usepackage{enumitem}
\usepackage{hyperref}
\usepackage{longtable}
\usepackage{setspace}
\usepackage{amssymb}

\definecolor{navy}{RGB}{21,42,79}
\definecolor{accent}{RGB}{37,99,235}

\titleformat{\chapter}[display]
  {\normalfont\centering\bfseries\color{navy}}
  {\large\MakeUppercase{\chaptertitlename}\ \thechapter}
  {4pt}
  {\Large}
\titlespacing*{\chapter}{0pt}{-8pt}{10pt}
\titleformat{\section}{\normalfont\large\bfseries\color{navy}}{\thesection}{8pt}{}
\titleformat{\subsection}{\normalfont\normalsize\bfseries}{\thesubsection}{6pt}{}

\setlength{\headheight}{14pt}
\pagestyle{fancy}
\fancyhf{}
\fancyhead[C]{\small\nouppercase{\leftmark}}
\fancyfoot[C]{\small\thepage}
\renewcommand{\headrulewidth}{0.4pt}
\fancypagestyle{plain}{\fancyhf{}\fancyhead[C]{\small\nouppercase{\leftmark}}\fancyfoot[C]{\small\thepage}\renewcommand{\headrulewidth}{0.4pt}}

\hypersetup{
    colorlinks=true,
    linkcolor=navy,
    citecolor=navy,
    urlcolor=accent,
    pdftitle={Weekly Assessment Report -- Used Car Price Prediction and Analysis System},
    pdfauthor={Bankar Smitraj Dinkar}
}

\setlength{\parindent}{0pt}
\setlength{\parskip}{5pt}
\setstretch{1.05}

% Full-border table helper (week metadata block; detailed work follows as itemize below)
% Attendance is intentionally not repeated here -- see the Consolidated
% Attendance Record chapter for the single authoritative attendance log.
\newcommand{\weektable}[3]{%
\begin{longtable}{|p{3.6cm}|p{10.5cm}|}
\hline
\textbf{Week} & \textbf{#1} \\
\hline
\textbf{Duration} & #2 \\
\hline
\textbf{Objectives} & #3 \\
\hline
\end{longtable}
}

\begin{document}
\pagenumbering{roman}

% =========================================================
% TITLE PAGE
% =========================================================
\begin{titlepage}
\centering
\thispagestyle{empty}
\vspace*{0.1cm}
\includegraphics[width=2.1cm]{images/college_logo.png}\\[4pt]
{\large\bfseries Sanjivani College of Engineering, Kopargaon\par}
{\normalsize Department of Computer Engineering\par}
{\normalsize (Affiliated to Savitribai Phule Pune University, Pune)\par}
\vspace{0.2cm}
\rule{\linewidth}{0.8pt}

\vspace{0.5cm}
{\Huge\bfseries\color{navy} Weekly Assessment Report\par}
\vspace{0.3cm}
{\Large\bfseries Used Car Price Prediction and Analysis System\\ Using Machine Learning\par}
\vspace{0.3cm}
{\large A Week-wise Progress and Assessment Log Submitted in Partial\\
Fulfilment of the Requirements for the Course\par}
\vspace{0.15cm}
{\Large\bfseries HOCO492 --- Mini Project: Artificial Intelligence and Machine Learning\par}

\vspace{0.65cm}
{\large\itshape Submitted by\par}
\vspace{0.15cm}
{\LARGE\bfseries Bankar Smitraj Dinkar\par}
{\large Roll No.: 09 \quad\textbar\quad PRN No.: UCS23M1009\par}

\vspace{0.65cm}
{\large\itshape Under the Guidance of\par}
\vspace{0.15cm}
{\Large\bfseries Prof. Bhagyashri Bhagawan Kotame\par}
{\large Assistant Professor, Department of Computer Engineering\par}

\vfill
{\large Department of Computer Engineering\par}
{\large Sanjivani College of Engineering, Kopargaon\par}
{\large Academic Year 2026--2027 \quad\textbar\quad Duration: 09-Jul-2026 to 18-Oct-2026 (15 Weeks)\par}
\end{titlepage}

% =========================================================
\chapter*{About This Report}
\thispagestyle{empty}

This document records the week-wise progress of the mini project
\textbf{``Used Car Price Prediction and Analysis System Using Machine
Learning''} over the full \textbf{15-week} project duration (09 July 2026
-- 18 October 2026), as required for continuous mentor assessment under
course \textbf{HOCO492}.

Each week lists the planned objectives and the actual work performed,
with references to the concrete files/artifacts produced in the project
repository.

\textbf{Note on structure:} the technical build of this project follows
the standard ML lifecycle (data acquisition $\to$ preprocessing $\to$ EDA
$\to$ modelling $\to$ evaluation $\to$ deployment $\to$ testing $\to$
documentation). That lifecycle is mapped below onto the 15 available
calendar weeks in the order the work was actually planned and completed.

\tableofcontents

\clearpage
\pagenumbering{arabic}
\setcounter{page}{1}

% =========================================================
\chapter{Week 1: Requirement Analysis and Topic Finalization}
\weektable{1}{06 Jul 2026 -- 12 Jul 2026}{Finalize the mini-project topic, define the problem statement, and align it with course outcomes CO1--CO4.}

\begin{itemize}[nosep]
    \item Selected and finalized the project topic: \textit{Used Car Price Prediction and Analysis System Using Machine Learning}.
    \item Defined the problem statement: used-car pricing is inconsistent and hard to estimate without domain expertise; a data-driven regression-based estimator is needed.
    \item Drafted project objectives and mapped them explicitly to course outcomes CO1--CO4 (independent investigation, orderly reporting, literature linkage, practical constraints).
    \item Reviewed the HOCO492 mini-project evaluation guidelines and scope (classical ML only, no deep learning).
\end{itemize}
\textbf{Deliverable:} \path{docs/project_proposal.md} (problem statement, objectives, scope, CO alignment).

% =========================================================
\chapter{Week 2: Literature Survey}
\weektable{2}{13 Jul 2026 -- 19 Jul 2026}{Survey published work on used-car / vehicle price prediction using machine learning and identify methodological gaps to address.}

\begin{itemize}[nosep]
    \item Reviewed and verified five real published studies: Pudaruth (2014), Monburinon et al. (2018, ICBIR), Gegic et al. (2019, \textit{TEM Journal}), Pal et al. (2018, FICC/Springer), and Noor \& Jan (2017, \textit{IJCA}).
    \item Cross-checked each paper's authors, venue, year, method, and reported result via independent search before citing.
    \item Synthesized a common finding across papers: tree-based ensembles consistently outperform plain linear regression for used-car pricing.
    \item Identified a methodological gap to avoid: several papers report only training accuracy without a comparable held-out check (motivated this project's explicit train/CV/test separation).
\end{itemize}
\textbf{Deliverable:} \path{docs/literature_survey.md} (5 verified references with per-paper relevance analysis).

% =========================================================
\chapter{Week 3: Dataset Acquisition and Project Setup}
\weektable{3}{20 Jul 2026 -- 26 Jul 2026}{Identify a real, public dataset suitable for the problem and set up the project's software structure and environment.}

\begin{itemize}[nosep]
    \item Located and validated the CarDekho used-car dataset (\textit{``Car details v3.csv''}, Kaggle \path{nehalbirla/vehicle-dataset-from-cardekho}) -- 8{,}128 rows, 13 columns, matching the required feature set.
    \item Set up the full project repository structure: \path{data/}, \path{src/}, \path{notebooks/}, \path{models/}, \path{reports/}, \path{tests/}, \path{docs/}.
    \item Configured the Python 3.13 environment with pandas, NumPy, scikit-learn, XGBoost, matplotlib, seaborn, Plotly, Streamlit, joblib, and pytest.
    \item Implemented \path{src/config.py} (paths, column mapping, reproducibility constants) and \path{src/data_loader.py} (schema-validated CSV loading).
\end{itemize}
\textbf{Deliverable:} \path{data/raw/used_cars.csv}, \path{src/config.py}, \path{src/data_loader.py}.

% =========================================================
\chapter{Week 4: Data Cleaning and Preprocessing Pipeline}
\weektable{4}{27 Jul 2026 -- 02 Aug 2026}{Clean the raw dataset and build a leakage-free preprocessing pipeline.}

\begin{itemize}[nosep]
    \item Implemented \path{clean_raw_data()}: duplicate removal (1{,}202 exact duplicates found), regex-based unit stripping (\texttt{"1248 CC"} $\to$ \texttt{1248.0}, \texttt{"74 bhp"} $\to$ \texttt{74.0}, \texttt{"23.4 kmpl"} $\to$ \texttt{23.4}), invalid-target row removal, and sanity filters on \path{km_driven}/\texttt{year}.
    \item Dropped the \texttt{torque} column after determining its free-text format could not be parsed reliably.
    \item Built the \texttt{ColumnTransformer}-based preprocessing pipeline: median imputation + standard scaling (numeric), most-frequent imputation + one-hot encoding (categorical).
\end{itemize}
\textbf{Deliverable:} \path{src/preprocessing.py}.

% =========================================================
\chapter{Week 5: Exploratory Data Analysis and Feature Engineering}
\weektable{5}{03 Aug 2026 -- 09 Aug 2026}{Explore the cleaned dataset and engineer features that improve model learnability.}

\begin{itemize}[nosep]
    \item Built and executed \path{notebooks/01_eda.ipynb} covering all required EDA items: distributions, missing-value and duplicate analysis, correlation matrix, price vs.\ age, km driven, mileage, and engine capacity, price by fuel type, transmission, and seller type, and outlier analysis.
    \item Implemented \path{car_age} (current year $-$ manufacturing year), \texttt{brand} extraction from the free-text \texttt{name} column, and rare-brand grouping ($<$20 listings $\to$ ``Other'').
    \item Saved all generated EDA plots to \path{reports/figures/}.
\end{itemize}
\textbf{Deliverable:} \path{notebooks/01_eda.ipynb}, \path{src/feature_engineering.py}, \path{reports/figures/*.png}.

% =========================================================
\chapter{Week 6: Baseline Model Development}
\weektable{6}{10 Aug 2026 -- 16 Aug 2026}{Implement the core training pipeline and establish baseline regression performance.}

\begin{itemize}[nosep]
    \item Implemented \path{src/train.py}: data preparation, 80/20 train--test split (\path{random_state=42}), and the candidate-model registry.
    \item Trained and evaluated baseline Linear Regression, Ridge Regression, and Decision Tree models.
    \item Established the baseline test R\textsuperscript{2} $\approx$ 0.71 (linear models) as the bar subsequent ensemble models must clear.
\end{itemize}
\textbf{Deliverable:} \path{src/train.py} (baseline stage), \path{notebooks/02_preprocessing.ipynb}.

% =========================================================
\chapter{Week 7: Ensemble Models and Cross-Validation}
\weektable{7}{17 Aug 2026 -- 23 Aug 2026}{Extend the candidate set to ensemble methods and compare all candidates via 5-fold cross-validation.}

\begin{itemize}[nosep]
    \item Added Random Forest, Gradient Boosting, and XGBoost to the candidate model set (6 algorithms total).
    \item Implemented 5-fold cross-validation (\path{cross_validate_models()}) reporting CV RMSE, R\textsuperscript{2}, and MAE (mean $\pm$ std) for every candidate on the training split only.
    \item Observed tree-based ensembles clearly outperforming linear models (CV R\textsuperscript{2} $\approx$ 0.87--0.88 vs. $\approx$ 0.69), consistent with the literature survey findings (Week 2).
\end{itemize}
\textbf{Deliverable:} \path{reports/cv_results.csv}.

% =========================================================
\chapter{Week 8: Hyperparameter Tuning}
\weektable{8}{24 Aug 2026 -- 30 Aug 2026}{Tune the strongest candidate models within a bounded search budget.}

\begin{itemize}[nosep]
    \item Defined bounded hyperparameter search spaces for Random Forest and XGBoost (n\_estimators, max\_depth, learning\_rate, subsample, etc.).
    \item Ran \texttt{RandomizedSearchCV} (20 iterations, 5-fold CV) on both models via \path{tune_model()}.
    \item Logged tuned CV RMSE and best hyperparameters for each; observed that tuning gains were modest and did not uniformly beat the untuned models on the later test-set check (Week 9).
\end{itemize}
\textbf{Deliverable:} Tuned model pipelines and logged best-parameter sets (see \path{src/train.py::tune_model}).

% =========================================================
\chapter{Week 9: Final Evaluation and Model Selection}
\weektable{9}{31 Aug 2026 -- 06 Sep 2026}{Evaluate every candidate once on the held-out test set and objectively select the final model.}

\begin{itemize}[nosep]
    \item Evaluated all 8 candidates (6 baseline + 2 tuned) exactly once on the untouched 20\% test split.
    \item Selected \textbf{Gradient Boosting} as the final model -- objectively the lowest test RMSE (Rs.\,151{,}043, R\textsuperscript{2} = 0.9281, MAPE = 21.28\%), narrowly ahead of the tuned Random Forest and XGBoost.
    \item Verified no evidence of overfitting (test R\textsuperscript{2} not lower than the CV R\textsuperscript{2} range).
    \item Generated diagnostic plots: actual-vs-predicted, residual plot, error distribution, feature importance.
    \item Serialized the final bundled pipeline (preprocessing + model) with \texttt{joblib}.
\end{itemize}
\textbf{Deliverable:} \path{reports/model_comparison.csv}, \path{reports/results.md}, \path{models/used_car_price_model.pkl}, \path{models/model_metadata.json}, \path{reports/figures/*.png}.

% =========================================================
\chapter{Week 10: Streamlit Application Development}
\weektable{10}{07 Sep 2026 -- 13 Sep 2026}{Build the interactive web application around the trained pipeline.}

\begin{itemize}[nosep]
    \item Implemented \path{src/predict.py}: input validation rules (year range, non-negative km, plausible mileage/engine/power bounds, known categorical values) and the prediction interface, including the empirical price-range calculation from test-set residual quantiles.
    \item Built \path{app.py} with five pages: \textit{Home, Price Prediction, Data Analysis, Model Performance, About Project} -- all reading live from generated reports/artifacts, nothing hard-coded.
    \item Manually verified the Price Prediction form end-to-end with sample vehicle inputs.
\end{itemize}
\textbf{Deliverable:} \path{app.py}, \path{src/predict.py}.

% =========================================================
\chapter{Week 11: Automated Testing and Reliability}
\weektable{11}{14 Sep 2026 -- 20 Sep 2026}{Validate the system with automated tests and benchmark prediction reliability.}

\begin{itemize}[nosep]
    \item Wrote 40 automated \texttt{pytest} test cases across \path{tests/test_data.py}, \path{test_preprocessing.py}, \path{test_model.py}, and \path{test_prediction.py} -- all 40 passing.
    \item Ran reliability testing: 50 repeated predictions on identical input (0\% variance, deterministic output) and measured per-prediction latency (mean 2.99\,ms on a warm model).
    \item Verified application deployment, routing, and UI responsiveness locally across all application views.
\end{itemize}
\textbf{Deliverable:} \path{tests/*.py}, \path{reports/test_cases.md}, \path{reports/reliability_test.md}.

% =========================================================
\chapter{Week 12: Project Documentation}
\weektable{12}{21 Sep 2026 -- 27 Sep 2026}{Produce the full supporting documentation set expected for the course (SDLC, design, methodology, limitations).}

\begin{itemize}[nosep]
    \item Authored \path{docs/sdlc.md} (Agile sprint plan mapped to this project's actual phases).
    \item Authored \path{docs/system_design.md} (layered architecture, data-flow tables) and \path{docs/methodology.md} (full pipeline, leakage-avoidance rationale).
    \item Authored \path{docs/limitations.md} and \path{docs/future_scope.md}, documenting honest constraints (regional scope, dropped \texttt{torque} feature, empirical price range) rather than overstating capability.
    \item Authored \path{docs/academic_report.md} (18-chapter report skeleton) and the top-level \path{README.md}.
\end{itemize}
\textbf{Deliverable:} \path{docs/*.md}, \path{README.md}.

% =========================================================
\chapter{Week 13: Project Synopsis Report (LaTeX)}
\weektable{13}{28 Sep 2026 -- 04 Oct 2026}{Prepare the formal, submission-ready project synopsis document.}

\begin{itemize}[nosep]
    \item Authored the full LaTeX synopsis (\path{reports/synopsis.tex}) following standard blackbook formatting: Times 12pt, justified body text, centered chapter headings, roman/arabic page numbering, footer pagination, List of Figures/Tables.
    \item Captured real application screenshots via automated browser testing (Playwright driving the live Streamlit app) for the Price Prediction, Data Analysis, and Model Performance pages.
    \item Added system architecture diagrams (layered architecture, end-to-end ML pipeline), the mathematical formulation of every algorithm and metric used, and full training-algorithm pseudocode.
    \item Compiled with MiKTeX/\texttt{pdflatex}; verified via page-by-page ink-coverage audit that no page is more than half blank.
\end{itemize}
\textbf{Deliverable:} \path{reports/synopsis.pdf}.

% =========================================================
\chapter{Week 14: Final Verification and Refinement}
\weektable{14}{05 Oct 2026 -- 11 Oct 2026}{Re-verify the entire system end-to-end and polish report formatting.}

\begin{itemize}[nosep]
    \item Re-ran \texttt{python -m src.train} end-to-end and confirmed all generated metrics/plots are reproducible from the included dataset with the fixed random seed.
    \item Re-ran the full \texttt{pytest} suite (40/40 passing) and re-verified the application runs and serves correctly.
    \item Refined report formatting per review feedback: fully bordered tables, justified/centered layout, footer-based pagination, proper font standards.
\end{itemize}
\textbf{Deliverable:} Finalized \path{reports/synopsis.pdf}; verified reproducibility end-to-end.

% =========================================================
\chapter{Week 15: Viva Preparation and Final Submission}
\weektable{15}{12 Oct 2026 -- 18 Oct 2026}{Prepare for the viva-voce examination and complete final project submission.}

\begin{itemize}[nosep]
    \item Prepared \path{docs/viva_questions.md} (44 anticipated questions with answers spanning ML fundamentals, preprocessing, evaluation metrics, and system design).
    \item Prepared \path{docs/presentation_outline.md} (12-slide defense presentation).
    \item Performed a final walkthrough of the live application and all generated reports as rehearsal for the viva demonstration.
    \item Assembled the complete project repository (code, trained model, tests, documentation, and this report) for final submission.
\end{itemize}
\textbf{Deliverable:} Complete, submission-ready project repository.

% =========================================================
\chapter*{Overall Completion Summary and Sign-off}
\addcontentsline{toc}{chapter}{Overall Completion Summary and Sign-off}

\section*{Milestone Completion Checklist}
\begin{longtable}{|p{1.2cm}|p{8.3cm}|p{3.2cm}|}
\hline
\textbf{Week} & \textbf{Milestone} & \textbf{Status} \\
\hline
1--2 & Topic finalization and literature survey & Completed \\
\hline
3--4 & Dataset acquisition, cleaning and preprocessing pipeline & Completed \\
\hline
5 & Exploratory data analysis and feature engineering & Completed \\
\hline
6--8 & Model development, cross-validation and hyperparameter tuning & Completed \\
\hline
9 & Final model evaluation and selection (Gradient Boosting) & Completed \\
\hline
10 & Streamlit application development & Completed \\
\hline
11 & Automated testing, reliability testing and containerization & Completed \\
\hline
12--13 & Documentation and project synopsis report & Completed \\
\hline
14 & Final verification and refinement & Completed \\
\hline
15 & Viva preparation and final submission & Completed \\
\hline
\end{longtable}

\section*{Overall Remarks}
\noindent\rule{\linewidth}{0.4pt}
\vspace{1.3cm}
\noindent\rule{\linewidth}{0.4pt}
\vspace{1.3cm}
\noindent\rule{\linewidth}{0.4pt}

\vspace{1cm}
\begin{minipage}[t]{0.3\linewidth}
\underline{\hspace{3.5cm}}\\
Student Signature\\
Bankar Smitraj Dinkar
\end{minipage}
\hfill
\begin{minipage}[t]{0.3\linewidth}
\underline{\hspace{3.5cm}}\\
Guide Signature\\
Prof. Bhagyashri B. Kotame
\end{minipage}
\hfill
\begin{minipage}[t]{0.3\linewidth}
\underline{\hspace{3.5cm}}\\
HOD Signature\\
Dr. M. A. Jawale
\end{minipage}

\end{document}
